\section*{Resumo}
\label{sec:resumo}

A deteção de fraude financeira sob extremo desequilíbrio de classes é um problema de aprendizagem automática com consequências económicas e operacionais diretas. Apesar do elevado volume de trabalho publicado, a literatura apresenta limitações metodológicas recorrentes — fuga de dados através de reamostragem incorreta, utilização de ROC-AUC como métrica principal em cenários de desequilíbrio severo, seleção implícita de limiar de decisão e avaliações em conjunto de dados único — que impedem comparações fiáveis entre famílias de modelos clássicos e de aprendizagem profunda.

Esta dissertação apresenta um estudo comparativo controlado e livre de fuga de dados envolvendo seis famílias de modelos — Regressão Logística, Random Forest, LightGBM, CatBoost, One-Class SVM e o FT-Transformer — em dois conjuntos de dados de referência: o ULB Credit Card Fraud e o Bank Account Fraud (BAF). Um desenho fatorial completo avalia sete estratégias de tratamento de desequilíbrio de classes em cinco modelos supervisionados e dois conjuntos de dados. A seleção de limiar de decisão é tratada como variável experimental de primeira ordem, com quatro estratégias comparadas. A generalização entre domínios é avaliada transferindo modelos treinados no BAF Base para cinco variantes com desvio de distribuição, sem reajuste. A interpretabilidade é analisada com SHAP e os padrões de atenção do FT-Transformer são examinados através de um enquadramento diagnóstico baseado em entropia.

Sob o protocolo livre de fuga de dados, os modelos de gradiente boosting superam o FT-Transformer no conjunto de dados ULB, de menor dimensão (CatBoost PR-AUC 0,885 vs FT-Transformer 0,856), enquanto ambos atingem desempenho idêntico no BAF, de maior escala (PR-AUC 0,180). A seleção de limiar revela-se crítica: fixar $\tau=0{,}5$ reduz $F_2$ a valores próximos de zero no BAF para todos os modelos, enquanto a estratégia max-$F_2$ recupera aproximadamente nove vezes esse valor. O SMOTE degrada o FT-Transformer em 41\% no BAF, comparado com apenas 14\% no CatBoost, revelando uma sensibilidade à sobreamostragem sintética específica à arquitetura. O Random Forest é o modelo mais robusto na generalização entre domínios, melhorando em quatro das cinco variantes do BAF apesar de apresentar o menor desempenho no domínio de treino. Os diagnósticos de atenção do FT-Transformer revelam um padrão de Visão Nebulosa (entropia normalizada 0,976), confirmando que o teto de desempenho de PR-AUC 0,18 reflete o limite informativo do espaço de características e não uma falha do modelo.

Os resultados fornecem um enquadramento de decisão para profissionais e uma explicação mecanística para a razão pela qual os modelos de gradiente boosting permanecem competitivos face a transformers tabulares sofisticados na deteção de fraude financeira.

\section*{Palavras-Chave}
\label{sec:palavras}

Deteção de fraude financeira, aprendizagem automática, gradient boosted decision trees, transformers tabulares, desequilíbrio de classes, avaliação livre de fuga de dados, SHAP, diagnósticos de atenção, generalização entre domínios.
